\section{From Diagnosis to Prompt: Echoing}
\label{sec:echoing}

The diagnosis prescribes the prompt with no design freedom left. If the answer
position cannot read the question, put something after the image that it can
read. There are exactly two things available.

\paragraph{Question echoing (\stit{}).} Repeat the question after the image.
This restores question-answer adjacency while provably leaving perception
untouched, since causal masking means a later copy cannot alter earlier patches
($\cos(\sti{}, \stit{}) = 0.99999$, Section~\ref{sec:anatomy-steering}). It
costs $13$ tokens, and it is the controlled intervention the diagnosis calls
for: matched on content, matched on perception, differing only in whether the
answer position has a question to read.

\paragraph{Image echoing (\sitit{}).} Repeat the image after the question, so
the answer position has both a nearby question and image tokens encoded with the
question already in context. Prior work also reversed the echoed copy
(\sititrev{}) and found it changed little, which we reproduce: NaturalBench
group accuracy moves from $37.37$ to $37.37$ ($p = 1.00$) and Winoground from
$40.25$ to $41.00$ ($p = 0.68$).

Two results in Table~\ref{tab:ladder} keep us honest. Question echoing alone is
not reliably better than simply asking the question last (NaturalBench $35.00$
against $35.05$, $p = 1.00$), so its role is to repair question-first prompting,
not to improve on question-last. Image echoing is a real gain and is the claim
we defend ($+0.085$ on Winoground, $p = 2.4\times10^{-5}$), though on POPE it
does not help at all.

What prior work left open is what the echo is made of and what it costs. It
listed the second image's token cost as a limitation and did not measure it.
Sections~\ref{sec:results} and~\ref{sec:ablations} take up that thread: how far
the paradox reaches, where the repair stops working, how little of the echo is
needed, and whether a prompt optimizer would have found it anyway.
